\section{Learning Objectives}

After completing this chapter, readers will be able to:
\begin{itemize}
    \item Apply a systematic methodology for designing complete end-to-end ML systems
    \item Evaluate and select appropriate architectures using decision frameworks
    \item Implement integration patterns that connect data, training, and serving pipelines
    \item Design systems that meet performance and scalability requirements
    \item Plan for system evolution, maintenance, and continuous improvement
    \item Analyze trade-offs between different architectural approaches
    \item Design production-ready ML systems from requirements to deployment
\end{itemize}

\section{Introduction}

Building end-to-end ML systems requires integrating concepts from all previous chapters: data engineering, model development, deployment, monitoring, and operations. This chapter provides a comprehensive methodology for designing, implementing, and evolving production ML systems that deliver business value while remaining maintainable and scalable.

Unlike previous chapters that focused on specific components, this chapter takes a holistic view of ML system architecture, examining how components interact and how to make informed trade-offs between competing requirements.

\textbf{Key Challenges in End-to-End ML Systems:}

\begin{itemize}
    \item \textbf{Complexity:} ML systems involve many interdependent components
    \item \textbf{Trade-offs:} Balancing latency, accuracy, cost, and maintainability
    \item \textbf{Evolution:} Systems must adapt to changing requirements and data
    \item \textbf{Integration:} Connecting ML components with existing infrastructure
    \item \textbf{Operationalization:} Moving from prototype to production-ready system
\end{itemize}

\section{System Design Methodology}

\subsection{The ML System Design Process}

A systematic approach to ML system design follows these phases:

\begin{enumerate}
    \item \textbf{Problem Framing:} Define business objectives and ML formulation
    \item \textbf{Requirements Analysis:} Gather technical and business constraints
    \item \textbf{Architecture Design:} Select components and integration patterns
    \item \textbf{Implementation Planning:} Break down into deliverable milestones
    \item \textbf{Validation Strategy:} Define success metrics and testing approach
    \item \textbf{Evolution Roadmap:} Plan for future enhancements and scaling
\end{enumerate}

\subsection{Phase 1: Problem Framing}

\textbf{Business Objectives Translation:}

\begin{table}[h]
\centering
\caption{Business Objective to ML Problem Mapping}
\begin{tabular}{p{4cm}p{4cm}p{5cm}}
\toprule
\textbf{Business Goal} & \textbf{ML Formulation} & \textbf{Success Metrics} \\
\midrule
Reduce customer churn & Binary classification & Precision, recall, F1-score, business value \\
\midrule
Personalize recommendations & Ranking/retrieval & Click-through rate, conversion, revenue lift \\
\midrule
Optimize pricing & Regression/RL & Revenue, margin, customer satisfaction \\
\midrule
Detect anomalies & Anomaly detection & Detection rate, false positive rate, MTTD \\
\midrule
Automate support & NLP classification & Resolution rate, user satisfaction, cost savings \\
\bottomrule
\end{tabular}
\end{table}

\textbf{Critical Questions for Problem Framing:}

\begin{itemize}
    \item What specific decision will the ML system enable or automate?
    \item What is the baseline (non-ML) performance?
    \item What improvement justifies the investment in ML?
    \item Are there regulatory or ethical constraints?
    \item What is the cost of false positives vs. false negatives?
    \item How will humans interact with the system?
\end{itemize}

\subsection{Phase 2: Requirements Analysis}

\textbf{Functional Requirements:}

\begin{lstlisting}[language=Python, caption=Requirements Specification Template]
#!/usr/bin/env python3
"""
ML System Requirements Specification
"""

from dataclasses import dataclass
from typing import List, Optional
from enum import Enum

class LatencyClass(Enum):
    REAL_TIME = "< 100ms"
    NEAR_REAL_TIME = "< 1s"
    BATCH = "> 1s"

class ConsistencyLevel(Enum):
    STRONG = "immediate consistency"
    EVENTUAL = "eventual consistency"

@dataclass
class FunctionalRequirements:
    """Functional requirements for ML system."""

    # Core capabilities
    prediction_types: List[str]  # e.g., ["classification", "ranking"]
    input_modalities: List[str]  # e.g., ["text", "images", "tabular"]
    output_format: str  # e.g., "probability", "ranking", "class"

    # Business rules
    business_constraints: List[str]
    regulatory_requirements: List[str]
    explainability_needed: bool
    human_in_the_loop: bool

@dataclass
class NonFunctionalRequirements:
    """Non-functional requirements for ML system."""

    # Performance
    latency_requirement: LatencyClass
    throughput_rps: int  # requests per second
    availability_target: float  # e.g., 0.999 for 99.9%

    # Scalability
    expected_users: int
    expected_growth_rate: float  # annual growth
    peak_to_average_ratio: float

    # Data
    training_data_size: str  # e.g., "10GB", "1TB"
    data_update_frequency: str  # e.g., "hourly", "daily"
    data_retention_period: str

    # Model
    model_accuracy_target: float
    model_update_frequency: str
    max_model_size: str  # e.g., "1GB"

    # Operations
    max_deployment_time: str  # e.g., "15 minutes"
    monitoring_granularity: str  # e.g., "per-request", "aggregate"
    incident_response_time: str  # e.g., "< 1 hour"

    # Cost
    infrastructure_budget: Optional[float]
    cost_per_prediction: Optional[float]

# Example instantiation
requirements = NonFunctionalRequirements(
    latency_requirement=LatencyClass.REAL_TIME,
    throughput_rps=10000,
    availability_target=0.999,
    expected_users=5_000_000,
    expected_growth_rate=0.5,  # 50% annual growth
    peak_to_average_ratio=3.0,
    training_data_size="500GB",
    data_update_frequency="daily",
    data_retention_period="2 years",
    model_accuracy_target=0.95,
    model_update_frequency="weekly",
    max_model_size="2GB",
    max_deployment_time="10 minutes",
    monitoring_granularity="aggregate",
    incident_response_time="30 minutes",
    infrastructure_budget=50000.0,  # monthly
    cost_per_prediction=0.001
)
\end{lstlisting}

\subsection{Phase 3: Architecture Design}

\textbf{Layered Architecture Pattern:}

\begin{figure}[h]
\centering
\resizebox{\textwidth}{!}{%
\begin{tikzpicture}[
    node distance=1.2cm,
    layer/.style={rectangle, draw, fill=blue!20, text width=10cm, text centered, minimum height=1cm},
    component/.style={rectangle, draw, fill=green!20, text width=2.3cm, text centered, minimum height=0.7cm, font=\footnotesize}
]

% Layers
\node[layer] (presentation) {Presentation Layer};
\node[layer, below=of presentation] (application) {Application Layer};
\node[layer, below=of application] (ml) {ML Layer};
\node[layer, below=of ml] (data) {Data Layer};
\node[layer, below=of data] (infrastructure) {Infrastructure Layer};

% Presentation components
\node[component, above right=0.3cm and -4.8cm of presentation] (web) {Web UI};
\node[component, right=0.1cm of web] (mobile) {Mobile App};
\node[component, right=0.1cm of mobile] (api) {REST API};
\node[component, right=0.1cm of api] (sdk) {SDK};

% Application components
\node[component, above right=0.3cm and -4.8cm of application] (auth) {Auth};
\node[component, right=0.1cm of auth] (business) {Business Logic};
\node[component, right=0.1cm of business] (orchestration) {Orchestration};
\node[component, right=0.1cm of orchestration] (cache) {Caching};

% ML components
\node[component, above right=0.3cm and -4.8cm of ml] (serving) {Model Serving};
\node[component, right=0.1cm of serving] (features) {Feature Store};
\node[component, right=0.1cm of features] (training) {Training};
\node[component, right=0.1cm of training] (registry) {Registry};

% Data components
\node[component, above right=0.3cm and -4.8cm of data] (streaming) {Streaming};
\node[component, right=0.1cm of streaming] (batch) {Batch ETL};
\node[component, right=0.1cm of batch] (warehouse) {Warehouse};
\node[component, right=0.1cm of warehouse] (lake) {Data Lake};

% Infrastructure components
\node[component, above right=0.3cm and -4.8cm of infrastructure] (compute) {Compute};
\node[component, right=0.1cm of compute] (storage) {Storage};
\node[component, right=0.1cm of storage] (network) {Network};
\node[component, right=0.1cm of network] (monitoring) {Monitoring};
\end{tikzpicture}}
\caption{Layered Architecture for End-to-End ML Systems}
\end{figure}

\section{Architecture Decision Frameworks}

\subsection{Decision Matrix for Architecture Patterns}

When designing ML systems, architects must choose between competing patterns. Use this framework to evaluate options systematically.

\textbf{Architecture Decision Record (ADR) Template:}

\begin{lstlisting}[language=Python, caption=ADR Implementation]
#!/usr/bin/env python3
"""
Architecture Decision Record (ADR) System
"""

from dataclasses import dataclass, field
from datetime import datetime
from typing import List, Dict
from enum import Enum

class DecisionStatus(Enum):
    PROPOSED = "proposed"
    ACCEPTED = "accepted"
    DEPRECATED = "deprecated"
    SUPERSEDED = "superseded"

@dataclass
class Alternative:
    """Alternative architectural approach."""
    name: str
    description: str
    pros: List[str]
    cons: List[str]
    cost_estimate: float
    complexity_score: int  # 1-10
    risk_score: int  # 1-10

@dataclass
class ArchitectureDecision:
    """Architecture Decision Record."""

    # Metadata
    id: int
    title: str
    date: datetime
    status: DecisionStatus
    decision_makers: List[str]

    # Context
    context: str  # Why is this decision needed?
    business_drivers: List[str]
    technical_drivers: List[str]
    constraints: List[str]

    # Analysis
    alternatives: List[Alternative]
    evaluation_criteria: Dict[str, float]  # criterion -> weight

    # Decision
    chosen_alternative: str
    rationale: str
    consequences: List[str]

    # Follow-up
    validation_metrics: List[str]
    review_date: datetime
    related_decisions: List[int] = field(default_factory=list)

class DecisionFramework:
    """Framework for evaluating architectural decisions."""

    def __init__(self):
        self.decisions: List[ArchitectureDecision] = []

    def score_alternative(
        self,
        alternative: Alternative,
        criteria: Dict[str, float]
    ) -> float:
        """
        Score an alternative based on weighted criteria.

        Args:
            alternative: Alternative to score
            criteria: Dict of criterion -> weight

        Returns:
            Weighted score (0-100)
        """
        scores = {
            'cost': (10 - alternative.cost_estimate / 10000) * 10,
            'complexity': (10 - alternative.complexity_score) * 10,
            'risk': (10 - alternative.risk_score) * 10,
            'pros': len(alternative.pros) * 10,
            'cons': (5 - len(alternative.cons)) * 10
        }

        weighted_score = sum(
            scores.get(criterion, 0) * weight
            for criterion, weight in criteria.items()
        )

        return weighted_score

    def recommend_alternative(
        self,
        alternatives: List[Alternative],
        criteria: Dict[str, float]
    ) -> Alternative:
        """Recommend best alternative based on criteria."""
        scored = [
            (alt, self.score_alternative(alt, criteria))
            for alt in alternatives
        ]
        return max(scored, key=lambda x: x[1])[0]

# Example: Batch vs Stream Processing Decision
batch_alternative = Alternative(
    name="Batch Processing",
    description="Daily batch jobs for feature generation",
    pros=[
        "Simpler to implement and debug",
        "Lower operational complexity",
        "Better for large-scale data processing",
        "Easier to reprocess historical data"
    ],
    cons=[
        "Higher latency (24-hour delay)",
        "Cannot handle real-time use cases",
        "Resource spikes during batch windows"
    ],
    cost_estimate=5000,  # monthly
    complexity_score=3,
    risk_score=2
)

stream_alternative = Alternative(
    name="Stream Processing",
    description="Real-time streaming with Kafka/Flink",
    pros=[
        "Low latency (seconds)",
        "Enables real-time features",
        "Smooth resource utilization",
        "Better user experience"
    ],
    cons=[
        "Higher operational complexity",
        "More difficult to debug",
        "Requires specialized expertise",
        "Harder to reprocess data"
    ],
    cost_estimate=15000,  # monthly
    complexity_score=8,
    risk_score=6
)

hybrid_alternative = Alternative(
    name="Hybrid Batch + Stream",
    description="Stream for real-time, batch for complex features",
    pros=[
        "Balanced latency and complexity",
        "Leverages strengths of both",
        "Can handle diverse use cases",
        "Incremental migration path"
    ],
    cons=[
        "Highest complexity overall",
        "Need to maintain two systems",
        "Potential consistency issues",
        "Higher total cost"
    ],
    cost_estimate=18000,  # monthly
    complexity_score=9,
    risk_score=7
)

# Create decision
decision = ArchitectureDecision(
    id=1,
    title="Data Processing Architecture",
    date=datetime.now(),
    status=DecisionStatus.PROPOSED,
    decision_makers=["ML Lead", "Data Engineering Lead", "CTO"],
    context="""
    Need to design data processing pipeline for real-time
    recommendation system with 10K RPS requirement.
    """,
    business_drivers=[
        "Reduce latency to improve conversion rates",
        "Enable personalization features",
        "Support A/B testing"
    ],
    technical_drivers=[
        "Handle 10K RPS",
        "Process 100GB daily data",
        "Support both batch and real-time features"
    ],
    constraints=[
        "Budget: $20K/month",
        "Team size: 3 data engineers",
        "Timeline: 3 months"
    ],
    alternatives=[batch_alternative, stream_alternative, hybrid_alternative],
    evaluation_criteria={
        'cost': 0.2,
        'complexity': 0.3,
        'risk': 0.2,
        'pros': 0.2,
        'cons': 0.1
    },
    chosen_alternative="Hybrid Batch + Stream",
    rationale="""
    While hybrid approach has highest complexity, it provides
    the best balance for our use case. We can start with batch
    processing for complex features and add streaming for
    real-time features incrementally.
    """,
    consequences=[
        "Need to hire streaming expert",
        "6-month timeline instead of 3",
        "Invest in monitoring and tooling"
    ],
    validation_metrics=[
        "Feature freshness < 1 minute",
        "System uptime > 99.9%",
        "Cost within budget"
    ],
    review_date=datetime(2025, 12, 31)
)
\end{lstlisting}

\subsection{Trade-off Analysis Framework}

\textbf{Common Trade-offs in ML Systems:}

\begin{table}[h]
\centering
\caption{Architecture Trade-off Matrix}
\begin{tabular}{p{3cm}p{3cm}p{3cm}p{3cm}}
\toprule
\textbf{Dimension} & \textbf{Option A} & \textbf{Option B} & \textbf{Consider When} \\
\midrule
Latency & Batch (high) & Real-time (low) & User-facing vs. backend \\
\midrule
Accuracy & Complex model & Simple model & Business value of improvement \\
\midrule
Cost & GPU training & CPU training & Budget vs. time-to-market \\
\midrule
Complexity & Monolith & Microservices & Team size and expertise \\
\midrule
Freshness & Static model & Continuous learning & Data drift rate \\
\midrule
Consistency & Strong & Eventual & Correctness requirements \\
\midrule
Coupling & Tight & Loose & System stability needs \\
\bottomrule
\end{tabular}
\end{table}

\section{Integration Patterns and Best Practices}

\subsection{Data Integration Patterns}

\textbf{Pattern 1: Event-Driven Architecture}

\begin{lstlisting}[language=Python, caption=Event-Driven ML System Integration]
#!/usr/bin/env python3
"""
Event-Driven Integration for ML Systems
"""

import json
from typing import Dict, List, Callable, Any
from dataclasses import dataclass, asdict
from datetime import datetime
import asyncio
from abc import ABC, abstractmethod

@dataclass
class Event:
    """Base event class."""
    event_id: str
    event_type: str
    timestamp: datetime
    source: str
    data: Dict[str, Any]

    def to_json(self) -> str:
        """Serialize event to JSON."""
        return json.dumps({
            **asdict(self),
            'timestamp': self.timestamp.isoformat()
        })

class EventBus(ABC):
    """Abstract event bus interface."""

    @abstractmethod
    async def publish(self, topic: str, event: Event):
        """Publish event to topic."""
        pass

    @abstractmethod
    async def subscribe(self, topic: str, handler: Callable):
        """Subscribe to topic with handler."""
        pass

class KafkaEventBus(EventBus):
    """Kafka-based event bus implementation."""

    def __init__(self, bootstrap_servers: List[str]):
        self.bootstrap_servers = bootstrap_servers
        self.handlers: Dict[str, List[Callable]] = {}

    async def publish(self, topic: str, event: Event):
        """Publish event to Kafka topic."""
        # Simplified - use aiokafka in production
        print(f"Publishing to {topic}: {event.to_json()}")

    async def subscribe(self, topic: str, handler: Callable):
        """Subscribe handler to topic."""
        if topic not in self.handlers:
            self.handlers[topic] = []
        self.handlers[topic].append(handler)

class MLSystemIntegration:
    """
    Integration layer for ML system components.
    Coordinates between data, training, and serving.
    """

    def __init__(self, event_bus: EventBus):
        self.event_bus = event_bus
        self._setup_handlers()

    def _setup_handlers(self):
        """Register event handlers for ML pipeline."""
        # Data pipeline events
        asyncio.create_task(
            self.event_bus.subscribe(
                'data.ingested',
                self.handle_data_ingestion
            )
        )

        # Training events
        asyncio.create_task(
            self.event_bus.subscribe(
                'training.completed',
                self.handle_training_completion
            )
        )

        # Monitoring events
        asyncio.create_task(
            self.event_bus.subscribe(
                'model.performance_degraded',
                self.handle_performance_degradation
            )
        )

    async def handle_data_ingestion(self, event: Event):
        """
        Handle new data ingestion.
        Triggers feature computation and model evaluation.
        """
        print(f"Data ingested: {event.data['record_count']} records")

        # Trigger feature computation
        feature_event = Event(
            event_id=f"feat-{event.event_id}",
            event_type="feature.compute_requested",
            timestamp=datetime.now(),
            source="integration_layer",
            data={
                'data_batch_id': event.data['batch_id'],
                'feature_set': 'v1'
            }
        )
        await self.event_bus.publish('feature.requests', feature_event)

        # Check if retraining needed
        if self._should_retrain(event):
            training_event = Event(
                event_id=f"train-{event.event_id}",
                event_type="training.requested",
                timestamp=datetime.now(),
                source="integration_layer",
                data={
                    'trigger': 'scheduled',
                    'data_version': event.data['data_version']
                }
            )
            await self.event_bus.publish('training.requests', training_event)

    async def handle_training_completion(self, event: Event):
        """
        Handle training completion.
        Validates model and triggers deployment if passing.
        """
        model_id = event.data['model_id']
        metrics = event.data['metrics']

        print(f"Training completed: {model_id}")
        print(f"Metrics: {metrics}")

        # Validate model
        if self._validate_model(metrics):
            # Trigger deployment
            deploy_event = Event(
                event_id=f"deploy-{event.event_id}",
                event_type="deployment.requested",
                timestamp=datetime.now(),
                source="integration_layer",
                data={
                    'model_id': model_id,
                    'strategy': 'canary',
                    'canary_percentage': 10
                }
            )
            await self.event_bus.publish('deployment.requests', deploy_event)
        else:
            # Alert on failed validation
            alert_event = Event(
                event_id=f"alert-{event.event_id}",
                event_type="alert.model_validation_failed",
                timestamp=datetime.now(),
                source="integration_layer",
                data={
                    'model_id': model_id,
                    'metrics': metrics,
                    'severity': 'high'
                }
            )
            await self.event_bus.publish('alerts', alert_event)

    async def handle_performance_degradation(self, event: Event):
        """
        Handle model performance degradation.
        Triggers retraining or rollback.
        """
        model_id = event.data['model_id']
        metric = event.data['metric']
        current_value = event.data['current_value']
        threshold = event.data['threshold']

        print(f"Performance degraded: {model_id}")
        print(f"{metric}: {current_value} (threshold: {threshold})")

        # Trigger retraining
        training_event = Event(
            event_id=f"train-{event.event_id}",
            event_type="training.requested",
            timestamp=datetime.now(),
            source="integration_layer",
            data={
                'trigger': 'performance_degradation',
                'model_to_replace': model_id
            }
        )
        await self.event_bus.publish('training.requests', training_event)

    def _should_retrain(self, event: Event) -> bool:
        """Determine if retraining is needed."""
        # Example logic
        return event.data.get('record_count', 0) > 1000000

    def _validate_model(self, metrics: Dict[str, float]) -> bool:
        """Validate model metrics meet thresholds."""
        return (
            metrics.get('accuracy', 0) > 0.90 and
            metrics.get('precision', 0) > 0.85 and
            metrics.get('recall', 0) > 0.85
        )

# Example usage
async def main():
    # Initialize event bus
    event_bus = KafkaEventBus(['localhost:9092'])

    # Initialize integration layer
    integration = MLSystemIntegration(event_bus)

    # Simulate data ingestion event
    data_event = Event(
        event_id="data-001",
        event_type="data.ingested",
        timestamp=datetime.now(),
        source="data_pipeline",
        data={
            'batch_id': 'batch-2025-001',
            'record_count': 1500000,
            'data_version': 'v2.1'
        }
    )

    await event_bus.publish('data.ingested', data_event)

    # Simulate training completion
    training_event = Event(
        event_id="train-001",
        event_type="training.completed",
        timestamp=datetime.now(),
        source="training_pipeline",
        data={
            'model_id': 'model-v2.1',
            'metrics': {
                'accuracy': 0.92,
                'precision': 0.89,
                'recall': 0.87
            }
        }
    )

    await event_bus.publish('training.completed', training_event)

# Run example
# asyncio.run(main())
\end{lstlisting}

\textbf{Pattern 2: API Gateway Integration}

\begin{lstlisting}[language=Python, caption=API Gateway for ML System]
#!/usr/bin/env python3
"""
API Gateway for ML System Integration
"""

from fastapi import FastAPI, HTTPException, Depends, Header
from fastapi.middleware.cors import CORSMiddleware
from pydantic import BaseModel, Field
from typing import Dict, List, Optional, Any
import time
import logging
from datetime import datetime
import asyncio
from circuitbreaker import circuit
import redis

# Configure logging
logging.basicConfig(level=logging.INFO)
logger = logging.getLogger(__name__)

# Models
class PredictionRequest(BaseModel):
    """Prediction request schema."""
    user_id: str
    context: Dict[str, Any]
    features: Optional[Dict[str, float]] = None
    model_version: Optional[str] = None

class PredictionResponse(BaseModel):
    """Prediction response schema."""
    prediction: Any
    confidence: float
    model_version: str
    latency_ms: float
    request_id: str

class HealthResponse(BaseModel):
    """Health check response."""
    status: str
    timestamp: datetime
    components: Dict[str, str]

# API Gateway
class MLAPIGateway:
    """
    API Gateway for ML system.
    Handles routing, load balancing, caching, and monitoring.
    """

    def __init__(
        self,
        model_endpoints: Dict[str, str],
        feature_store_url: str,
        cache_url: str
    ):
        self.model_endpoints = model_endpoints
        self.feature_store_url = feature_store_url
        self.cache = redis.Redis.from_url(cache_url)

        # Metrics
        self.request_count = 0
        self.error_count = 0
        self.total_latency = 0.0

    async def predict(
        self,
        request: PredictionRequest,
        model_name: str = "default"
    ) -> PredictionResponse:
        """
        Handle prediction request.
        Implements caching, feature enrichment, and error handling.
        """
        start_time = time.time()
        request_id = f"req-{int(time.time())}-{self.request_count}"
        self.request_count += 1

        try:
            # Check cache
            cache_key = self._get_cache_key(request, model_name)
            cached = self._get_from_cache(cache_key)
            if cached:
                logger.info(f"Cache hit: {request_id}")
                return cached

            # Enrich with features from feature store
            if request.features is None:
                request.features = await self._get_features(request.user_id)

            # Get prediction from model service
            prediction = await self._call_model_service(
                model_name,
                request
            )

            # Calculate latency
            latency_ms = (time.time() - start_time) * 1000
            self.total_latency += latency_ms

            # Build response
            response = PredictionResponse(
                prediction=prediction['prediction'],
                confidence=prediction['confidence'],
                model_version=prediction['model_version'],
                latency_ms=latency_ms,
                request_id=request_id
            )

            # Cache response
            self._cache_response(cache_key, response)

            # Log metrics
            logger.info(
                f"Request {request_id}: "
                f"latency={latency_ms:.2f}ms, "
                f"confidence={prediction['confidence']:.3f}"
            )

            return response

        except Exception as e:
            self.error_count += 1
            logger.error(f"Prediction error: {e}")
            raise HTTPException(status_code=500, detail=str(e))

    @circuit(failure_threshold=5, recovery_timeout=60)
    async def _call_model_service(
        self,
        model_name: str,
        request: PredictionRequest
    ) -> Dict[str, Any]:
        """
        Call model service with circuit breaker.
        Falls back to backup model if primary fails.
        """
        endpoint = self.model_endpoints.get(model_name)
        if not endpoint:
            raise ValueError(f"Unknown model: {model_name}")

        try:
            # Simulate model service call
            await asyncio.sleep(0.01)  # Simulated latency

            # Return prediction
            return {
                'prediction': 0.85,
                'confidence': 0.92,
                'model_version': 'v1.2.3'
            }
        except Exception as e:
            logger.warning(f"Primary model failed: {e}, trying backup")
            return await self._call_backup_model(request)

    async def _call_backup_model(
        self,
        request: PredictionRequest
    ) -> Dict[str, Any]:
        """Fallback to backup model."""
        # Simplified backup model
        return {
            'prediction': 0.5,  # Default prediction
            'confidence': 0.5,
            'model_version': 'backup-v1.0'
        }

    async def _get_features(self, user_id: str) -> Dict[str, float]:
        """Fetch features from feature store."""
        # Simulate feature store call
        await asyncio.sleep(0.005)

        return {
            'feature_1': 0.5,
            'feature_2': 0.8,
            'feature_3': 0.3
        }

    def _get_cache_key(
        self,
        request: PredictionRequest,
        model_name: str
    ) -> str:
        """Generate cache key."""
        return f"pred:{model_name}:{request.user_id}:{hash(str(request.context))}"

    def _get_from_cache(self, key: str) -> Optional[PredictionResponse]:
        """Get prediction from cache."""
        try:
            cached = self.cache.get(key)
            if cached:
                return PredictionResponse.parse_raw(cached)
        except:
            pass
        return None

    def _cache_response(
        self,
        key: str,
        response: PredictionResponse,
        ttl: int = 300
    ):
        """Cache prediction response."""
        try:
            self.cache.setex(
                key,
                ttl,
                response.json()
            )
        except Exception as e:
            logger.warning(f"Cache write failed: {e}")

    def get_metrics(self) -> Dict[str, Any]:
        """Get gateway metrics."""
        avg_latency = (
            self.total_latency / self.request_count
            if self.request_count > 0 else 0
        )
        error_rate = (
            self.error_count / self.request_count
            if self.request_count > 0 else 0
        )

        return {
            'request_count': self.request_count,
            'error_count': self.error_count,
            'error_rate': error_rate,
            'avg_latency_ms': avg_latency
        }

# FastAPI application
app = FastAPI(title="ML API Gateway")

# Add CORS
app.add_middleware(
    CORSMiddleware,
    allow_origins=["*"],
    allow_methods=["*"],
    allow_headers=["*"],
)

# Initialize gateway
gateway = MLAPIGateway(
    model_endpoints={
        'default': 'http://model-service:8000',
        'recommendation': 'http://rec-service:8001',
        'fraud': 'http://fraud-service:8002'
    },
    feature_store_url='http://feature-store:8003',
    cache_url='redis://localhost:6379'
)

# Authentication dependency
async def verify_api_key(x_api_key: str = Header(...)):
    """Verify API key."""
    if x_api_key != "secret-key":  # Use proper auth in production
        raise HTTPException(status_code=401, detail="Invalid API key")
    return x_api_key

# Endpoints
@app.post("/predict/{model_name}", response_model=PredictionResponse)
async def predict(
    model_name: str,
    request: PredictionRequest,
    api_key: str = Depends(verify_api_key)
):
    """Prediction endpoint."""
    return await gateway.predict(request, model_name)

@app.get("/health", response_model=HealthResponse)
async def health_check():
    """Health check endpoint."""
    return HealthResponse(
        status="healthy",
        timestamp=datetime.now(),
        components={
            'api_gateway': 'up',
            'cache': 'up',
            'model_service': 'up'
        }
    )

@app.get("/metrics")
async def metrics():
    """Metrics endpoint."""
    return gateway.get_metrics()
\end{lstlisting}

\subsection{Model Integration Patterns}

\textbf{Pattern 3: Model Ensemble Integration}

\begin{lstlisting}[language=Python, caption=Ensemble Model Integration]
#!/usr/bin/env python3
"""
Ensemble Model Integration Pattern
"""

from typing import List, Dict, Any, Optional
from dataclasses import dataclass
import numpy as np
from abc import ABC, abstractmethod
import asyncio

@dataclass
class ModelPrediction:
    """Single model prediction."""
    model_id: str
    prediction: Any
    confidence: float
    latency_ms: float
    metadata: Dict[str, Any]

class EnsembleStrategy(ABC):
    """Abstract ensemble strategy."""

    @abstractmethod
    def combine(
        self,
        predictions: List[ModelPrediction]
    ) -> ModelPrediction:
        """Combine predictions from multiple models."""
        pass

class AveragingEnsemble(EnsembleStrategy):
    """Averaging ensemble strategy."""

    def combine(
        self,
        predictions: List[ModelPrediction]
    ) -> ModelPrediction:
        """Average predictions."""
        pred_values = [p.prediction for p in predictions]
        confidences = [p.confidence for p in predictions]

        return ModelPrediction(
            model_id="ensemble_average",
            prediction=np.mean(pred_values),
            confidence=np.mean(confidences),
            latency_ms=max(p.latency_ms for p in predictions),
            metadata={'n_models': len(predictions)}
        )

class WeightedEnsemble(EnsembleStrategy):
    """Weighted ensemble strategy."""

    def __init__(self, weights: Dict[str, float]):
        self.weights = weights

    def combine(
        self,
        predictions: List[ModelPrediction]
    ) -> ModelPrediction:
        """Weighted average of predictions."""
        weighted_sum = sum(
            p.prediction * self.weights.get(p.model_id, 1.0)
            for p in predictions
        )
        total_weight = sum(
            self.weights.get(p.model_id, 1.0)
            for p in predictions
        )

        return ModelPrediction(
            model_id="ensemble_weighted",
            prediction=weighted_sum / total_weight,
            confidence=max(p.confidence for p in predictions),
            latency_ms=max(p.latency_ms for p in predictions),
            metadata={
                'n_models': len(predictions),
                'weights': self.weights
            }
        )

class StackingEnsemble(EnsembleStrategy):
    """Stacking ensemble with meta-learner."""

    def __init__(self, meta_model):
        self.meta_model = meta_model

    def combine(
        self,
        predictions: List[ModelPrediction]
    ) -> ModelPrediction:
        """Use meta-model to combine predictions."""
        # Stack predictions as features for meta-model
        features = np.array([
            [p.prediction, p.confidence]
            for p in predictions
        ]).flatten()

        # Meta-model prediction
        meta_pred = self.meta_model.predict(features.reshape(1, -1))[0]

        return ModelPrediction(
            model_id="ensemble_stacking",
            prediction=meta_pred,
            confidence=0.9,  # From meta-model
            latency_ms=max(p.latency_ms for p in predictions) + 5,
            metadata={'n_models': len(predictions)}
        )

class EnsembleOrchestrator:
    """
    Orchestrates ensemble of models.
    Handles parallel execution and result combination.
    """

    def __init__(
        self,
        models: Dict[str, Any],
        strategy: EnsembleStrategy,
        timeout: float = 1.0
    ):
        self.models = models
        self.strategy = strategy
        self.timeout = timeout

    async def predict(
        self,
        features: Dict[str, Any]
    ) -> ModelPrediction:
        """
        Get predictions from all models and combine.

        Args:
            features: Input features

        Returns:
            Combined prediction
        """
        # Execute models in parallel
        tasks = [
            self._predict_single(model_id, model, features)
            for model_id, model in self.models.items()
        ]

        try:
            predictions = await asyncio.gather(
                *tasks,
                return_exceptions=True
            )

            # Filter out exceptions
            valid_predictions = [
                p for p in predictions
                if isinstance(p, ModelPrediction)
            ]

            if not valid_predictions:
                raise RuntimeError("All models failed")

            # Combine predictions
            return self.strategy.combine(valid_predictions)

        except asyncio.TimeoutError:
            # Fallback to fastest responding model
            return await self._fallback_prediction(features)

    async def _predict_single(
        self,
        model_id: str,
        model: Any,
        features: Dict[str, Any]
    ) -> ModelPrediction:
        """Get prediction from single model."""
        import time
        start = time.time()

        try:
            # Simulate model prediction
            await asyncio.sleep(0.01)
            prediction = 0.8  # model.predict(features)
            confidence = 0.9

            latency_ms = (time.time() - start) * 1000

            return ModelPrediction(
                model_id=model_id,
                prediction=prediction,
                confidence=confidence,
                latency_ms=latency_ms,
                metadata={}
            )
        except Exception as e:
            raise RuntimeError(f"Model {model_id} failed: {e}")

    async def _fallback_prediction(
        self,
        features: Dict[str, Any]
    ) -> ModelPrediction:
        """Fallback when ensemble fails."""
        return ModelPrediction(
            model_id="fallback",
            prediction=0.5,
            confidence=0.5,
            latency_ms=0.0,
            metadata={'reason': 'timeout'}
        )

# Example usage
async def main():
    # Mock models
    models = {
        'model_a': None,  # XGBoost
        'model_b': None,  # LightGBM
        'model_c': None   # Neural Network
    }

    # Weighted ensemble
    strategy = WeightedEnsemble({
        'model_a': 0.4,
        'model_b': 0.35,
        'model_c': 0.25
    })

    orchestrator = EnsembleOrchestrator(models, strategy)

    # Get prediction
    features = {'user_id': '123', 'item_id': '456'}
    result = await orchestrator.predict(features)

    print(f"Ensemble prediction: {result.prediction}")
    print(f"Confidence: {result.confidence}")
    print(f"Latency: {result.latency_ms}ms")

# asyncio.run(main())
\end{lstlisting}

\section{Performance and Scalability Considerations}

\subsection{Performance Optimization Strategies}

\textbf{1. Latency Optimization}

\begin{table}[h]
\centering
\caption{Latency Optimization Techniques}
\begin{tabular}{p{3.5cm}p{4cm}p{4cm}}
\toprule
\textbf{Technique} & \textbf{Latency Reduction} & \textbf{Trade-offs} \\
\midrule
Caching & 80-95\% (cache hits) & Stale data, memory cost \\
\midrule
Model quantization & 40-60\% & 1-2\% accuracy loss \\
\midrule
Batching & 30-50\% (throughput) & Increased latency per request \\
\midrule
Feature precomputation & 50-70\% & Storage cost, freshness \\
\midrule
Model pruning & 30-40\% & Slight accuracy loss \\
\midrule
GPU inference & 70-90\% (vs CPU) & Higher cost \\
\bottomrule
\end{tabular}
\end{table}

\textbf{2. Scalability Patterns}

\begin{lstlisting}[language=Python, caption=Auto-Scaling ML Service]
#!/usr/bin/env python3
"""
Auto-Scaling Configuration for ML Service
"""

from dataclasses import dataclass
from typing import List, Dict
from enum import Enum

class ScalingMetric(Enum):
    """Metrics for autoscaling decisions."""
    CPU_UTILIZATION = "cpu"
    MEMORY_UTILIZATION = "memory"
    REQUEST_COUNT = "requests"
    LATENCY_P99 = "latency_p99"
    QUEUE_DEPTH = "queue_depth"

@dataclass
class ScalingPolicy:
    """Auto-scaling policy configuration."""

    # Scale-up triggers
    scale_up_metric: ScalingMetric
    scale_up_threshold: float
    scale_up_evaluation_periods: int
    scale_up_cooldown_seconds: int

    # Scale-down triggers
    scale_down_metric: ScalingMetric
    scale_down_threshold: float
    scale_down_evaluation_periods: int
    scale_down_cooldown_seconds: int

    # Capacity limits
    min_instances: int
    max_instances: int
    target_instances: int

    # Scaling increments
    scale_up_increment: int = 2
    scale_down_increment: int = 1

@dataclass
class ResourceConfiguration:
    """Resource configuration for ML service."""

    # Compute
    cpu_cores: int
    memory_gb: int
    gpu_count: int
    gpu_type: str

    # Storage
    disk_gb: int

    # Network
    max_concurrent_requests: int
    request_timeout_seconds: int

class MLServiceScalingConfig:
    """Complete scaling configuration for ML service."""

    def __init__(self):
        # Resource configuration per instance
        self.resources = ResourceConfiguration(
            cpu_cores=4,
            memory_gb=16,
            gpu_count=1,
            gpu_type="T4",
            disk_gb=100,
            max_concurrent_requests=100,
            request_timeout_seconds=30
        )

        # Scaling policy
        self.scaling_policy = ScalingPolicy(
            scale_up_metric=ScalingMetric.CPU_UTILIZATION,
            scale_up_threshold=70.0,  # 70% CPU
            scale_up_evaluation_periods=2,  # 2 periods above threshold
            scale_up_cooldown_seconds=300,  # 5 minutes
            scale_down_metric=ScalingMetric.CPU_UTILIZATION,
            scale_down_threshold=30.0,  # 30% CPU
            scale_down_evaluation_periods=5,  # 5 periods below threshold
            scale_down_cooldown_seconds=600,  # 10 minutes
            min_instances=2,  # For high availability
            max_instances=20,
            target_instances=3,
            scale_up_increment=2,  # Scale up by 2 instances
            scale_down_increment=1  # Scale down by 1 instance
        )

    def get_kubernetes_manifest(self) -> Dict:
        """Generate Kubernetes HPA manifest."""
        return {
            'apiVersion': 'autoscaling/v2',
            'kind': 'HorizontalPodAutoscaler',
            'metadata': {
                'name': 'ml-service-hpa'
            },
            'spec': {
                'scaleTargetRef': {
                    'apiVersion': 'apps/v1',
                    'kind': 'Deployment',
                    'name': 'ml-service'
                },
                'minReplicas': self.scaling_policy.min_instances,
                'maxReplicas': self.scaling_policy.max_instances,
                'metrics': [
                    {
                        'type': 'Resource',
                        'resource': {
                            'name': 'cpu',
                            'target': {
                                'type': 'Utilization',
                                'averageUtilization': int(
                                    self.scaling_policy.scale_up_threshold
                                )
                            }
                        }
                    }
                ],
                'behavior': {
                    'scaleUp': {
                        'stabilizationWindowSeconds': (
                            self.scaling_policy.scale_up_cooldown_seconds
                        ),
                        'policies': [
                            {
                                'type': 'Pods',
                                'value': self.scaling_policy.scale_up_increment,
                                'periodSeconds': 60
                            }
                        ]
                    },
                    'scaleDown': {
                        'stabilizationWindowSeconds': (
                            self.scaling_policy.scale_down_cooldown_seconds
                        ),
                        'policies': [
                            {
                                'type': 'Pods',
                                'value': self.scaling_policy.scale_down_increment,
                                'periodSeconds': 60
                            }
                        ]
                    }
                }
            }
        }

# Example configuration
config = MLServiceScalingConfig()
hpa_manifest = config.get_kubernetes_manifest()
\end{lstlisting}

\textbf{3. Load Balancing Strategies}

\begin{itemize}
    \item \textbf{Round-robin:} Simple, even distribution (good for homogeneous instances)
    \item \textbf{Least-connections:} Route to instance with fewest connections
    \item \textbf{Weighted:} Distribute based on instance capacity
    \item \textbf{Latency-based:} Route to lowest-latency instance
    \item \textbf{Model-aware:} Route to specialized models based on request characteristics
\end{itemize}

\section{Evolution and Maintenance Strategies}

\subsection{System Evolution Framework}

\textbf{Evolution Maturity Levels:}

\begin{table}[h]
\centering
\caption{ML System Maturity Model}
\begin{tabular}{p{2cm}p{4cm}p{5cm}}
\toprule
\textbf{Level} & \textbf{Characteristics} & \textbf{Capabilities} \\
\midrule
Level 0: Manual & Scripts, notebooks & One-off model training \\
\midrule
Level 1: Automated Training & Training pipelines & Reproducible training \\
\midrule
Level 2: Automated Deployment & CI/CD for ML & Automated model deployment \\
\midrule
Level 3: Automated Monitoring & Monitoring + alerts & Detect degradation \\
\midrule
Level 4: Automated Retraining & Trigger-based retraining & Self-healing system \\
\midrule
Level 5: Full MLOps & End-to-end automation & Continuous learning \\
\bottomrule
\end{tabular}
\end{table}

\subsection{Maintenance Strategies}

\begin{lstlisting}[language=Python, caption=System Maintenance Framework]
#!/usr/bin/env python3
"""
ML System Maintenance Framework
"""

from dataclasses import dataclass
from datetime import datetime, timedelta
from typing import List, Dict, Optional
from enum import Enum

class MaintenanceType(Enum):
    """Types of maintenance activities."""
    PREVENTIVE = "preventive"
    CORRECTIVE = "corrective"
    ADAPTIVE = "adaptive"
    PERFECTIVE = "perfective"

@dataclass
class MaintenanceTask:
    """Maintenance task definition."""
    task_id: str
    task_type: MaintenanceType
    description: str
    frequency: str  # e.g., "daily", "weekly", "monthly"
    estimated_duration: timedelta
    impact: str  # "low", "medium", "high"
    requires_downtime: bool

class MLSystemMaintenance:
    """Framework for ML system maintenance."""

    def __init__(self):
        self.tasks = self._define_maintenance_tasks()

    def _define_maintenance_tasks(self) -> List[MaintenanceTask]:
        """Define all maintenance tasks."""
        return [
            # Preventive maintenance
            MaintenanceTask(
                task_id="prev-001",
                task_type=MaintenanceType.PREVENTIVE,
                description="Monitor data drift and feature distributions",
                frequency="daily",
                estimated_duration=timedelta(hours=1),
                impact="low",
                requires_downtime=False
            ),
            MaintenanceTask(
                task_id="prev-002",
                task_type=MaintenanceType.PREVENTIVE,
                description="Review model performance metrics",
                frequency="daily",
                estimated_duration=timedelta(hours=2),
                impact="low",
                requires_downtime=False
            ),
            MaintenanceTask(
                task_id="prev-003",
                task_type=MaintenanceType.PREVENTIVE,
                description="Check infrastructure health and resource usage",
                frequency="daily",
                estimated_duration=timedelta(minutes=30),
                impact="low",
                requires_downtime=False
            ),
            MaintenanceTask(
                task_id="prev-004",
                task_type=MaintenanceType.PREVENTIVE,
                description="Validate data pipeline integrity",
                frequency="weekly",
                estimated_duration=timedelta(hours=3),
                impact="medium",
                requires_downtime=False
            ),
            MaintenanceTask(
                task_id="prev-005",
                task_type=MaintenanceType.PREVENTIVE,
                description="Update dependencies and security patches",
                frequency="monthly",
                estimated_duration=timedelta(hours=4),
                impact="medium",
                requires_downtime=True
            ),

            # Corrective maintenance
            MaintenanceTask(
                task_id="corr-001",
                task_type=MaintenanceType.CORRECTIVE,
                description="Retrain model due to performance degradation",
                frequency="as_needed",
                estimated_duration=timedelta(hours=8),
                impact="high",
                requires_downtime=False
            ),
            MaintenanceTask(
                task_id="corr-002",
                task_type=MaintenanceType.CORRECTIVE,
                description="Fix data quality issues",
                frequency="as_needed",
                estimated_duration=timedelta(hours=4),
                impact="high",
                requires_downtime=False
            ),
            MaintenanceTask(
                task_id="corr-003",
                task_type=MaintenanceType.CORRECTIVE,
                description="Rollback failed deployment",
                frequency="as_needed",
                estimated_duration=timedelta(minutes=30),
                impact="critical",
                requires_downtime=True
            ),

            # Adaptive maintenance
            MaintenanceTask(
                task_id="adapt-001",
                task_type=MaintenanceType.ADAPTIVE,
                description="Add new features based on changing patterns",
                frequency="quarterly",
                estimated_duration=timedelta(days=5),
                impact="medium",
                requires_downtime=False
            ),
            MaintenanceTask(
                task_id="adapt-002",
                task_type=MaintenanceType.ADAPTIVE,
                description="Adapt to new data sources or formats",
                frequency="as_needed",
                estimated_duration=timedelta(days=3),
                impact="high",
                requires_downtime=False
            ),

            # Perfective maintenance
            MaintenanceTask(
                task_id="perf-001",
                task_type=MaintenanceType.PERFECTIVE,
                description="Optimize model latency",
                frequency="quarterly",
                estimated_duration=timedelta(days=7),
                impact="medium",
                requires_downtime=False
            ),
            MaintenanceTask(
                task_id="perf-002",
                task_type=MaintenanceType.PERFECTIVE,
                description="Improve model accuracy",
                frequency="quarterly",
                estimated_duration=timedelta(days=10),
                impact="medium",
                requires_downtime=False
            ),
            MaintenanceTask(
                task_id="perf-003",
                task_type=MaintenanceType.PERFECTIVE,
                description="Reduce infrastructure costs",
                frequency="quarterly",
                estimated_duration=timedelta(days=5),
                impact="low",
                requires_downtime=False
            )
        ]

    def get_maintenance_schedule(
        self,
        start_date: datetime,
        end_date: datetime
    ) -> List[tuple[datetime, MaintenanceTask]]:
        """
        Generate maintenance schedule for date range.

        Args:
            start_date: Schedule start date
            end_date: Schedule end date

        Returns:
            List of (scheduled_datetime, task) tuples
        """
        schedule = []

        for task in self.tasks:
            if task.frequency == "as_needed":
                continue

            # Calculate occurrences
            current_date = start_date
            while current_date <= end_date:
                schedule.append((current_date, task))

                # Increment based on frequency
                if task.frequency == "daily":
                    current_date += timedelta(days=1)
                elif task.frequency == "weekly":
                    current_date += timedelta(weeks=1)
                elif task.frequency == "monthly":
                    current_date += timedelta(days=30)
                elif task.frequency == "quarterly":
                    current_date += timedelta(days=90)
                else:
                    break

        return sorted(schedule, key=lambda x: x[0])

    def estimate_maintenance_effort(
        self,
        period_days: int = 30
    ) -> Dict[str, float]:
        """
        Estimate maintenance effort for period.

        Args:
            period_days: Period in days

        Returns:
            Dict with effort estimates by maintenance type
        """
        effort = {
            'preventive': 0.0,
            'corrective': 0.0,
            'adaptive': 0.0,
            'perfective': 0.0,
            'total': 0.0
        }

        for task in self.tasks:
            if task.frequency == "as_needed":
                continue

            # Calculate occurrences in period
            if task.frequency == "daily":
                occurrences = period_days
            elif task.frequency == "weekly":
                occurrences = period_days / 7
            elif task.frequency == "monthly":
                occurrences = period_days / 30
            elif task.frequency == "quarterly":
                occurrences = period_days / 90
            else:
                occurrences = 0

            # Add to effort
            hours = occurrences * task.estimated_duration.total_seconds() / 3600
            effort[task.task_type.value] += hours
            effort['total'] += hours

        return effort

# Example usage
maintenance = MLSystemMaintenance()

# Get 30-day schedule
start = datetime.now()
end = start + timedelta(days=30)
schedule = maintenance.get_maintenance_schedule(start, end)

print(f"Maintenance schedule ({len(schedule)} tasks):")
for date, task in schedule[:5]:
    print(f"  {date.strftime('%Y-%m-%d')}: {task.description}")

# Estimate effort
effort = maintenance.estimate_maintenance_effort(30)
print(f"\nEstimated 30-day effort: {effort['total']:.1f} hours")
print(f"  Preventive: {effort['preventive']:.1f}h")
print(f"  Corrective: {effort['corrective']:.1f}h")
print(f"  Adaptive: {effort['adaptive']:.1f}h")
print(f"  Perfective: {effort['perfective']:.1f}h")
\end{lstlisting}

\section{Real-World System Examples}

\subsection{Example 1: Real-Time Content Moderation System}

\textbf{Business Context:}
\begin{itemize}
    \item Social media platform with 50M daily active users
    \item Need to detect and remove harmful content (hate speech, violence, spam)
    \item Requirements: <100ms latency, 99.5\% recall, 95\% precision
    \item Handle multiple content types: text, images, video
\end{itemize}

\textbf{System Architecture:}

\begin{figure}[h]
\centering
\begin{tikzpicture}[
    node distance=1.3cm,
    component/.style={rectangle, draw, fill=blue!20, text width=2.5cm, text centered, minimum height=0.8cm, font=\footnotesize},
    storage/.style={cylinder, draw, fill=green!20, text width=2.2cm, text centered, minimum height=0.8cm, aspect=0.3, font=\footnotesize}
]

% User layer
\node[component] (user) {User Posts Content};

% API Gateway
\node[component, below=of user] (gateway) {API Gateway};

% Content processing
\node[component, below left=1cm and 0.5cm of gateway] (text) {Text Classifier};
\node[component, below=of gateway] (image) {Image Classifier};
\node[component, below right=1cm and 0.5cm of gateway] (video) {Video Analyzer};

% Aggregation
\node[component, below=2.5cm of gateway] (ensemble) {Ensemble Scorer};

% Decision
\node[component, below=of ensemble] (decision) {Decision Engine};

% Actions
\node[component, below left=1cm and 0.5cm of decision] (allow) {Allow Post};
\node[component, below=of decision] (review) {Human Review Queue};
\node[component, below right=1cm and 0.5cm of decision] (block) {Block Content};

% Storage
\node[storage, right=2cm of decision] (db) {Moderation DB};

% Arrows
\draw[->] (user) -- (gateway);
\draw[->] (gateway) -- (text);
\draw[->] (gateway) -- (image);
\draw[->] (gateway) -- (video);
\draw[->] (text) -- (ensemble);
\draw[->] (image) -- (ensemble);
\draw[->] (video) -- (ensemble);
\draw[->] (ensemble) -- (decision);
\draw[->] (decision) -- (allow);
\draw[->] (decision) -- (review);
\draw[->] (decision) -- (block);
\draw[->] (decision) -- (db);

\end{tikzpicture}
\caption{Content Moderation System Architecture}
\end{figure}

\textbf{Key Design Decisions:}

\begin{enumerate}
    \item \textbf{Parallel Processing:} Text, image, and video classifiers run in parallel to minimize latency
    \item \textbf{Ensemble Scoring:} Combine signals from multiple models for robust detection
    \item \textbf{Human-in-the-Loop:} Route borderline cases to human review (0.7 < score < 0.9)
    \item \textbf{Model Serving:} Use TorchServe with GPU inference for image/video models
    \item \textbf{Caching:} Cache recent decisions for duplicate content (40\% cache hit rate)
    \item \textbf{Monitoring:} Track precision/recall, false positive rate, and user appeals
\end{enumerate}

\textbf{Implementation Highlights:}

\begin{lstlisting}[language=Python, caption=Content Moderation Decision Engine]
#!/usr/bin/env python3
"""
Content Moderation Decision Engine
"""

from dataclasses import dataclass
from typing import List, Dict, Optional
from enum import Enum

class ContentType(Enum):
    TEXT = "text"
    IMAGE = "image"
    VIDEO = "video"

class ModerationAction(Enum):
    ALLOW = "allow"
    REVIEW = "review"
    BLOCK = "block"

@dataclass
class ModerationScore:
    """Moderation score from a single model."""
    content_type: ContentType
    violation_type: str  # "hate_speech", "violence", etc.
    score: float  # 0.0 (safe) to 1.0 (harmful)
    confidence: float

@dataclass
class ModerationDecision:
    """Final moderation decision."""
    action: ModerationAction
    combined_score: float
    reason: str
    model_scores: List[ModerationScore]
    requires_human_review: bool

class ModerationEngine:
    """
    Decision engine for content moderation.
    Combines signals from multiple models and applies business rules.
    """

    def __init__(
        self,
        block_threshold: float = 0.9,
        review_threshold: float = 0.7,
        weights: Optional[Dict[str, float]] = None
    ):
        self.block_threshold = block_threshold
        self.review_threshold = review_threshold
        self.weights = weights or {
            'hate_speech': 1.0,
            'violence': 1.0,
            'spam': 0.7,
            'nudity': 0.8
        }

    def make_decision(
        self,
        scores: List[ModerationScore]
    ) -> ModerationDecision:
        """
        Make moderation decision based on model scores.

        Args:
            scores: List of scores from different models

        Returns:
            Final moderation decision
        """
        # Calculate weighted combined score
        combined_score = self._combine_scores(scores)

        # Apply thresholds
        if combined_score >= self.block_threshold:
            action = ModerationAction.BLOCK
            reason = self._get_block_reason(scores)
            requires_review = False
        elif combined_score >= self.review_threshold:
            action = ModerationAction.REVIEW
            reason = "Borderline case requiring human judgment"
            requires_review = True
        else:
            action = ModerationAction.ALLOW
            reason = "Content passes moderation checks"
            requires_review = False

        # Check for high-confidence violations
        if self._has_high_confidence_violation(scores):
            action = ModerationAction.BLOCK
            requires_review = False

        return ModerationDecision(
            action=action,
            combined_score=combined_score,
            reason=reason,
            model_scores=scores,
            requires_human_review=requires_review
        )

    def _combine_scores(self, scores: List[ModerationScore]) -> float:
        """Combine scores using weighted average."""
        if not scores:
            return 0.0

        weighted_sum = sum(
            score.score * self.weights.get(score.violation_type, 1.0)
            for score in scores
        )
        total_weight = sum(
            self.weights.get(score.violation_type, 1.0)
            for score in scores
        )

        return weighted_sum / total_weight if total_weight > 0 else 0.0

    def _has_high_confidence_violation(
        self,
        scores: List[ModerationScore]
    ) -> bool:
        """Check if any model has high-confidence violation."""
        return any(
            score.score >= 0.95 and score.confidence >= 0.9
            for score in scores
        )

    def _get_block_reason(self, scores: List[ModerationScore]) -> str:
        """Get human-readable reason for blocking."""
        violations = [
            score.violation_type
            for score in scores
            if score.score >= self.review_threshold
        ]

        if len(violations) == 1:
            return f"Content violates policy: {violations[0]}"
        else:
            return f"Content violates multiple policies: {', '.join(violations)}"

# Example usage
engine = ModerationEngine()

scores = [
    ModerationScore(
        content_type=ContentType.TEXT,
        violation_type="hate_speech",
        score=0.92,
        confidence=0.88
    ),
    ModerationScore(
        content_type=ContentType.IMAGE,
        violation_type="violence",
        score=0.45,
        confidence=0.75
    )
]

decision = engine.make_decision(scores)
print(f"Action: {decision.action.value}")
print(f"Score: {decision.combined_score:.3f}")
print(f"Reason: {decision.reason}")
\end{lstlisting}

\textbf{Results:}
\begin{itemize}
    \item Achieved 99.7\% recall, 94\% precision
    \item Reduced manual review workload by 60\%
    \item p99 latency: 85ms
    \item Cost: \$0.0003 per moderation
\end{itemize}

\subsection{Example 2: Dynamic Pricing System}

\textbf{Business Context:}
\begin{itemize}
    \item Ride-sharing platform with 10M weekly trips
    \item Need to optimize prices based on supply/demand
    \item Requirements: Maximize revenue while maintaining acceptable wait times
    \item Handle surge pricing during peak hours and special events
\end{itemize}

\textbf{System Architecture:}

\begin{lstlisting}[language=Python, caption=Dynamic Pricing System]
#!/usr/bin/env python3
"""
Dynamic Pricing System for Ride-Sharing
"""

import numpy as np
from dataclasses import dataclass
from typing import Dict, List, Optional
from datetime import datetime, timedelta

@dataclass
class MarketConditions:
    """Current market conditions for pricing."""
    location_id: str
    available_drivers: int
    pending_requests: int
    avg_wait_time_minutes: float
    hour_of_day: int
    day_of_week: int
    is_special_event: bool
    weather_condition: str
    historical_demand: float

@dataclass
class PricingDecision:
    """Pricing decision output."""
    base_price: float
    surge_multiplier: float
    final_price: float
    expected_wait_time: float
    confidence: float
    reasoning: str

class DynamicPricingEngine:
    """
    ML-powered dynamic pricing engine.
    Uses demand forecasting, supply estimation, and optimization.
    """

    def __init__(
        self,
        base_price: float = 10.0,
        min_multiplier: float = 1.0,
        max_multiplier: float = 5.0
    ):
        self.base_price = base_price
        self.min_multiplier = min_multiplier
        self.max_multiplier = max_multiplier

    def calculate_price(
        self,
        conditions: MarketConditions
    ) -> PricingDecision:
        """
        Calculate optimal price based on market conditions.

        Args:
            conditions: Current market conditions

        Returns:
            Pricing decision
        """
        # Calculate supply-demand ratio
        supply_demand_ratio = (
            conditions.available_drivers /
            max(conditions.pending_requests, 1)
        )

        # Base surge multiplier on supply-demand
        if supply_demand_ratio >= 1.5:
            # Oversupply: no surge
            surge_multiplier = 1.0
        elif supply_demand_ratio >= 0.8:
            # Balanced: minimal surge
            surge_multiplier = 1.2
        elif supply_demand_ratio >= 0.5:
            # High demand: moderate surge
            surge_multiplier = 2.0
        else:
            # Very high demand: high surge
            surge_multiplier = 3.5

        # Adjust for time of day
        if conditions.hour_of_day in [7, 8, 9, 17, 18, 19]:
            # Peak commute hours
            surge_multiplier *= 1.2
        elif conditions.hour_of_day >= 22 or conditions.hour_of_day <= 5:
            # Late night/early morning
            surge_multiplier *= 1.3

        # Adjust for special events
        if conditions.is_special_event:
            surge_multiplier *= 1.5

        # Adjust for weather
        if conditions.weather_condition in ['rain', 'snow']:
            surge_multiplier *= 1.3

        # Cap multiplier
        surge_multiplier = np.clip(
            surge_multiplier,
            self.min_multiplier,
            self.max_multiplier
        )

        # Calculate final price
        final_price = self.base_price * surge_multiplier

        # Estimate wait time
        expected_wait_time = self._estimate_wait_time(
            conditions,
            surge_multiplier
        )

        # Generate reasoning
        reasoning = self._generate_reasoning(
            supply_demand_ratio,
            surge_multiplier,
            conditions
        )

        return PricingDecision(
            base_price=self.base_price,
            surge_multiplier=surge_multiplier,
            final_price=final_price,
            expected_wait_time=expected_wait_time,
            confidence=0.85,
            reasoning=reasoning
        )

    def _estimate_wait_time(
        self,
        conditions: MarketConditions,
        surge_multiplier: float
    ) -> float:
        """Estimate wait time based on conditions and pricing."""
        # Higher prices attract more drivers and reduce demand
        elasticity = 0.3  # 30% demand reduction per 2x price increase
        demand_adjustment = 1 - (elasticity * np.log2(surge_multiplier))

        adjusted_requests = conditions.pending_requests * demand_adjustment

        # Simple queuing model
        if conditions.available_drivers == 0:
            return 30.0  # Max wait time

        wait_time = (adjusted_requests / conditions.available_drivers) * 3.0
        return min(wait_time, 30.0)

    def _generate_reasoning(
        self,
        supply_demand_ratio: float,
        surge_multiplier: float,
        conditions: MarketConditions
    ) -> str:
        """Generate human-readable pricing reasoning."""
        reasons = []

        if supply_demand_ratio < 0.5:
            reasons.append("high demand")
        elif supply_demand_ratio > 1.5:
            reasons.append("low demand")

        if conditions.is_special_event:
            reasons.append("special event")

        if conditions.weather_condition in ['rain', 'snow']:
            reasons.append("adverse weather")

        if conditions.hour_of_day in [7, 8, 9, 17, 18, 19]:
            reasons.append("peak hours")

        if not reasons:
            return "Normal market conditions"

        return f"Surge pricing due to: {', '.join(reasons)}"

# Example usage
engine = DynamicPricingEngine()

# Peak hour with high demand
conditions = MarketConditions(
    location_id="downtown_01",
    available_drivers=15,
    pending_requests=50,
    avg_wait_time_minutes=8.5,
    hour_of_day=18,
    day_of_week=4,  # Friday
    is_special_event=False,
    weather_condition="clear",
    historical_demand=45.0
)

decision = engine.calculate_price(conditions)
print(f"Base price: ${decision.base_price:.2f}")
print(f"Surge multiplier: {decision.surge_multiplier:.2f}x")
print(f"Final price: ${decision.final_price:.2f}")
print(f"Expected wait: {decision.expected_wait_time:.1f} minutes")
print(f"Reasoning: {decision.reasoning}")
\end{lstlisting}

\textbf{Key Design Decisions:}

\begin{enumerate}
    \item \textbf{Real-Time Computation:} Price calculated on-demand based on current conditions
    \item \textbf{Demand Forecasting:} Use historical data + external signals (events, weather)
    \item \textbf{Price Elasticity:} Model how price changes affect demand
    \item \textbf{Fairness Constraints:} Cap maximum surge multiplier at 5x
    \item \textbf{Transparency:} Provide clear reasoning for surge pricing
    \item \textbf{A/B Testing:} Continuously test pricing strategies
\end{enumerate}

\textbf{Results:}
\begin{itemize}
    \item 23\% revenue increase vs. fixed pricing
    \item 18\% reduction in average wait times
    \item Driver earnings increased 15\%
    \item Customer satisfaction maintained above 4.2/5.0
\end{itemize}

\section{Chapter Summary}

Building end-to-end ML systems requires careful consideration of:

\begin{itemize}
    \item \textbf{System Design Methodology:} Follow structured process from problem framing to implementation
    \item \textbf{Architecture Decisions:} Use frameworks to evaluate trade-offs systematically
    \item \textbf{Integration Patterns:} Implement robust patterns for connecting components
    \item \textbf{Performance:} Design for scalability and low latency from the start
    \item \textbf{Evolution:} Plan for continuous improvement and adaptation
    \item \textbf{Maintenance:} Establish processes for ongoing system health
\end{itemize}

\textbf{Key Takeaways:}

\begin{bestpractice}
Start simple and iterate. Begin with minimal viable system, measure performance, and enhance based on real-world feedback. Premature optimization leads to complexity without proven value.
\end{bestpractice}

\begin{bestpractice}
Design for observability from day one. You cannot improve what you cannot measure. Instrument every component with logging, metrics, and tracing.
\end{bestpractice}

\begin{bestpractice}
Embrace trade-offs. Perfect systems don't exist. Make informed decisions about accuracy vs. latency, cost vs. performance, and complexity vs. maintainability.
\end{bestpractice}

\begin{bestpractice}
Plan for failure. Systems will fail. Design with redundancy, graceful degradation, and automated recovery mechanisms.
\end{bestpractice}

\section{Exercises}

\begin{exercise}
\textbf{Exercise 1: E-Commerce Search and Ranking System}

Design a complete end-to-end ML system for product search and ranking on an e-commerce platform.

\textbf{Requirements:}
\begin{itemize}
    \item 1 million products, 5 million monthly users
    \item Support text search with semantic understanding
    \item Personalized ranking based on user history
    \item Real-time inventory and pricing updates
    \item <200ms p99 latency
    \item A/B testing capability
\end{itemize}

\textbf{Deliverables:}
\begin{enumerate}
    \item Complete architecture diagram with all components
    \item Data pipeline design (batch and real-time)
    \item ML model architecture (retrieval + ranking)
    \item Serving infrastructure design
    \item Monitoring and evaluation strategy
    \item Cost estimate and capacity planning
    \item Discussion of key trade-offs and decisions
\end{enumerate}

\textbf{Consider:}
\begin{itemize}
    \item How will you handle cold-start for new users/products?
    \item How will you balance relevance, recency, and business metrics?
    \item What features will you use for ranking?
    \item How will you ensure diversity in results?
\end{itemize}
\end{exercise}

\begin{exercise}
\textbf{Exercise 2: Credit Risk Assessment System}

Design an ML system for assessing credit risk and determining loan approvals.

\textbf{Requirements:}
\begin{itemize}
    \item Process 10,000 loan applications daily
    \item Predict default probability with 90\%+ AUC
    \item Provide explanations for decisions (regulatory requirement)
    \item Fair lending compliance (no bias by protected attributes)
    \item Support manual override by loan officers
    \item Audit trail for all decisions
\end{itemize}

\textbf{Deliverables:}
\begin{enumerate}
    \item End-to-end system architecture
    \item Data sources and feature engineering strategy
    \item Model selection and training approach
    \item Explainability implementation
    \item Fairness evaluation and mitigation plan
    \item Deployment strategy and rollback plan
    \item Monitoring dashboard specification
\end{enumerate}

\textbf{Consider:}
\begin{itemize}
    \item How will you ensure model fairness across demographics?
    \item What explainability method will you use?
    \item How will you handle model updates without disrupting operations?
    \item What metrics will you monitor in production?
\end{itemize}
\end{exercise}

\begin{exercise}
\textbf{Exercise 3: Video Content Recommendation System}

Design a recommendation system for a video streaming platform similar to Netflix or YouTube.

\textbf{Requirements:}
\begin{itemize}
    \item 50 million users, 100,000 videos
    \item Personalized recommendations on homepage and autoplay
    \item Support cold-start for new users
    \item Optimize for watch time and user satisfaction
    \item Handle diverse content (movies, series, shorts)
    \item Support A/B testing and experimentation
\end{itemize}

\textbf{Deliverables:}
\begin{enumerate}
    \item Multi-stage recommendation architecture (candidate generation + ranking)
    \item Feature engineering for users and videos
    \item Model architecture (collaborative filtering, deep learning, etc.)
    \item Real-time serving infrastructure
    \item Evaluation metrics (online and offline)
    \item Experimentation framework
    \item Scalability plan for 10x growth
\end{enumerate}

\textbf{Consider:}
\begin{itemize}
    \item How will you balance exploration vs. exploitation?
    \item How will you prevent filter bubbles?
    \item What's your strategy for cold-start users?
    \item How will you handle content freshness?
\end{itemize}
\end{exercise}

\begin{exercise}
\textbf{Exercise 4: Predictive Maintenance System}

Design an ML system for predictive maintenance of industrial equipment.

\textbf{Requirements:}
\begin{itemize}
    \item Monitor 10,000 machines across 50 facilities
    \item Sensor data: temperature, vibration, pressure (sampled at 100Hz)
    \item Predict failures 7-14 days in advance
    \item Minimize false alarms (high cost of unnecessary maintenance)
    \item Provide actionable insights for maintenance teams
    \item Handle network connectivity issues at remote facilities
\end{itemize}

\textbf{Deliverables:}
\begin{enumerate}
    \item IoT data collection and processing architecture
    \item Feature engineering from time-series sensor data
    \item Anomaly detection and failure prediction models
    \item Edge vs. cloud processing strategy
    \item Alert generation and prioritization system
    \item Visualization dashboard for maintenance teams
    \item ROI analysis (cost savings from preventing failures)
\end{enumerate}

\textbf{Consider:}
\begin{itemize}
    \item How will you handle class imbalance (rare failures)?
    \item What features will you extract from time-series data?
    \item How will you handle missing data from sensor failures?
    \item What's your strategy for different equipment types?
\end{itemize}
\end{exercise}
